% !TeX root = ../../Book5.tex
% 纲要标题：### E.2 Picard–Vessiot 扩张与对应
% 纲要位置：工作记录/计划纲要.md，第 4695 行。
% * 在特征零、常数域代数闭的标准情形下介绍 PV 环、基本解矩阵、微分单性与无新常数条件
% * 保持基域且与导子交换的自同构构成常数域上的线性代数群；证明存在唯一性及群闭性
% * 中间微分域与 Zariski 闭子群的对应；闭性条件不能删去，给出证明概要
% #### E.2.1 PV 环、基本解矩阵与无新常数条件
% #### E.2.2 微分自同构群与存在唯一性
% #### E.2.3 中间微分域与 Zariski 闭子群
\section{Picard–Vessiot 扩张与对应}
\label{b5:sec:E02}

一个方程可能没有足够多的有理函数解。为研究全部线性解，应当加入基本解矩阵，同时控制新增的常数。这里始终固定特征零微分域 \((K,\delta)\)，并要求常数域 \(C=\ker\delta\) 代数闭。这个假设保证下述存在唯一性具有最简洁的形式。

\subsection{PV 环、基本解矩阵与无新常数条件}
\label{b5:subsec:E02:01}

\begin{definition}\label{b5:def:E-pv-ring}
设 \((K,\delta)\) 为特征零微分域，常数域 \(C=\ker\delta\) 代数闭，\(n\ge1\)，\(A\in M_n(K)\)。设交换微分 \(K\) 代数 \(R\) 是整环。若存在 \(Z\in\operatorname{GL}_n(R)\) 满足 \(\delta(Z)=AZ\)，有 \(R=K\Bigl[Z_{ij},(\det Z)^{-1}\Bigr]\)，且 \(R\) 的微分理想只有 \(0,R\)，则称 \(R\) 为系统 \(\delta(Y)=AY\) 的\term[Picard-Vessiot-huan]{Picard–Vessiot 环}[Picard--Vessiot ring]，简称 PV 环。这里微分理想指满足 \(\delta(I)\subset I\) 的理想；\(Z\) 称为\term[jiben-jie-juzhen]{基本解矩阵}[fundamental solution matrix]。分式域 \(L=\operatorname{Frac}R\) 称为相应的\term[Picard-Vessiot-kuozhang]{Picard–Vessiot 扩张}[Picard--Vessiot extension]。
\end{definition}

微分单性不等于环是域：稍后会出现 \(K[t,t^{-1}]\) 这样的 PV 环。其作用是排除与方程无关的微分商；而在上述常数假设下，PV 分式域满足 \(C_L=C\)，这称为无新常数。反过来，由基本解矩阵的系数生成、且无新常数的微分域扩张，也给出这样的 PV 环。下一小节证明这些性质。

\ProseParagraphBreak
若只要求存在解矩阵，就可以把 \(Z\) 乘以任意新增的超越常数，因而无法确定扩张。无新常数限制的是全部分式域，不能只检查所选解的几个系数。基本解矩阵也不是逐列任取非零解：必须要求行列式非零。

\subsection{微分自同构群与存在唯一性}
\label{b5:subsec:E02:02}

\begin{lemma}\label{b5:lem:E-simple-constants}
设 \((K,\delta)\) 为特征零微分域，常数域 \(C=\ker\delta\) 代数闭。若非零交换微分 \(K\) 代数 \(S\) 作为 \(K\) 代数有限生成，且其微分理想只有 \(0,S\)，则 \(S\) 是整环，且 \(\operatorname{Frac}S\) 的常数域为 \(C\)。
\end{lemma}
\begin{proof}
若 \(a\in S\) 不幂零，则
\(I_a=\{b:a^mb=0\text{ 对某个 }m\ge1\}\)
是一个真微分理想：对 \(a^mb=0\) 求导再乘 \(a\)，得到 \(a^{m+1}b'=0\)。故 \(I_a=0\)，\(a\) 不是零因子。若非零 \(a\) 幂零，取最小 \(m\) 使 \(a^m=0\)，则
\(ma^{m-1}a'=0\) 且 \(ma^{m-1}\ne0\)，所以 \(a'\) 是零因子，按刚才的结论必幂零。于是幂零元构成微分理想，只能为零。这证明 \(S\) 为整环。

若 \(c\in\operatorname{Frac}S\) 满足 \(c'=0\)，非零理想
\(\{b\in S:bc\in S\}\) 对导子稳定，故为 \(S\)，从而 \(c\in S\)。若 \(c\notin C\)，则对每个 \(a\in C\)，\((c-a)S\) 是非零微分理想，所以 \(c-a\) 都可逆。

这迫使 \(c\) 在 \(K\) 上代数。确实，若 \(c\) 超越，在 \(K(c)\) 上对 \(S\otimes_{K[c]}K(c)\) 作 Noether 规范化并清去有限多个分母，可取非零 \(f\in K[T]\)，使 \(S[f(c)^{-1}]\) 整于某个多项式环
\(K[c,f(c)^{-1}][u_1,\ldots,u_d]\)。取 \(a\in C\) 使 \(f(a)\ne0\)，整扩张的 lying-over 性质便给出一个包含 \(c-a\) 的素理想，与 \(c-a\) 可逆矛盾。最后，\(c\) 的首一最小多项式求导后，最高项消失；最小性说明其全部系数的导数为零。该多项式因此在 \(C[T]\) 中，\(C\) 代数闭又推出 \(c\in C\)。
\end{proof}

\begin{theorem}[Picard–Vessiot 存在唯一性]\label{b5:thm:E-pv-existence}
设 \((K,\delta)\) 为特征零微分域，常数域 \(C=\ker\delta\) 代数闭，\(n\ge1\)，\(A\in M_n(K)\)。则系统 \(\delta(Y)=AY\) 存在 PV 环，任意两个这样的环在微分 \(K\) 代数意义下同构；其分式域的常数域均为 \(C\)。反之，若 \(L/K\) 是微分域扩张，\(Z\in\operatorname{GL}_n(L)\) 满足 \(\delta(Z)=AZ\)，\(L=K(Z_{ij})\) 且 \(C_L=C\)，则 \(K\Bigl[Z_{ij},(\det Z)^{-1}\Bigr]\) 是该系统的 PV 环。
\end{theorem}

\begin{proof}
在 \(S_0=K[X_{ij},(\det X)^{-1}]\) 上按 \(X'=AX\) 延伸导子，取一个极大真微分理想 \(I\)。商环 \(R=S_0/I\) 微分单，且仍有可逆基本解矩阵。由\cref{b5:lem:E-simple-constants}，它是整环，分式域没有新常数，所以是所需 PV 环。

若 \(R_1,R_2\) 是两个 PV 环，在 \(R_1\otimes_KR_2\) 中取极大真微分理想并作商，得到微分单环 \(T\)。两个 \(R_i\to T\) 都单射，且由同一引理 \(T\) 的常数为 \(C\)。两基本解矩阵在 \(T\) 中的比值导数为零，故相差一个 \(\operatorname{GL}_n(C)\) 矩阵。它们的系数及逆行列式因而生成同一个子环，于是两个嵌入的像相等，给出所需同构。

为证反向陈述，需要一个常数变量引理：若微分域 \(F\) 的常数为 \(C\)，且变量 \(U_{ij}\) 的导数全为零，则
\(F[U_{ij},(\det U)^{-1}]\) 中的每个微分理想都由它与 \(C[U_{ij},(\det U)^{-1}]\) 的交生成。取常数多项式环的一组 \(C\) 基，把理想元写成有限项 \(\sum a_\alpha e_\alpha\)。按非零系数个数归纳：先除以一个非零系数，使其成为一；求导后这一项消失。若另一个系数 \(a\) 不为常数，对原式除以 \(a\) 再求导也减少一项，两式相减便给出 \(-(a'/a^2)\) 乘原式。归纳和 \(a'\ne0\) 说明原式由常数关系生成；全部系数已为常数时结论直接成立。

现设 \(L=K(Z_{ij})\)、\(Z'=AZ\)、\(C_L=C\)。仍取 \(S_0\) 的极大真微分理想 \(I\)，在
\(L\otimes_KS_0\) 中改变量为 \(X=ZU\)。此时 \(U'=0\)，刚才的引理把 \(I\) 的延伸化成由常数多项式理想 \(J\) 生成的理想。它是真理想，因为 \(L\otimes_K(S_0/I)\ne0\)。由弱 Nullstellensatz，\(J\) 有一个可逆常数矩阵零点 \(c\)。代入 \(X=Zc\) 得到微分同态 \(S_0\to L\)，其核包含 \(I\)；极大微分性使二者相等。其像正好是
\(K[Z_{ij},(\det Z)^{-1}]\)，故此环微分单，证明反向陈述。完整证明的相同代数构造见\cite[Section 1.3, Lemmas 1.17 and 1.23, Propositions 1.20 and 1.22]{VanDerPutSinger2003GaloisTheory}。
\end{proof}

同构一般不唯一，其不唯一性正是下一定义要研究的对象。概念及存在唯一性的概述见\cite[Definition 1.3.2 and the following discussion, pp. 10--11]{Singer2009DifferentialGalois}。

\begin{definition}\label{b5:def:E-galois-group}
设 \((K,\delta)\) 为特征零微分域，常数域 \(C=\ker\delta\) 代数闭，\(L/K\) 为 PV 扩张。所有固定 \(K\) 且满足 \(\sigma\delta=\delta\sigma\) 的 \(L\) 的域自同构组成群
\[
 G=\operatorname{Gal}_{\delta}(L/K).
\]
这个群称为该扩张的\term[weifen-Galois-qun]{微分 Galois 群}[differential Galois group]。
\end{definition}

\begin{proposition}\label{b5:prop:E-group-matrix}
设 \((K,\delta)\) 为特征零微分域，常数域 \(C=\ker\delta\) 代数闭，\(n\ge1\)，\(A\in M_n(K)\)。设 \(L/K\) 为系统 \(\delta(Y)=AY\) 的 PV 扩张，\(Z\in\operatorname{GL}_n(L)\) 为基本解矩阵。每个 \(\sigma\in\operatorname{Gal}_\delta(L/K)\) 唯一写成 \(\sigma(Z)=ZC_\sigma\)，其中 \(C_\sigma\in\operatorname{GL}_n(C)\)；对应 \(\sigma\mapsto C_\sigma\) 是忠实群表示。换用另一基本解矩阵只使此表示共轭。
\end{proposition}
\begin{proof}
利用 \(\delta(Z^{-1})=-Z^{-1}\delta(Z)Z^{-1}\)，直接求导得到 \(\delta\Bigl(Z^{-1}\sigma(Z)\Bigr)=0\)。无新常数保证该矩阵在 \(C\) 上。因为自同构固定常数，\(\sigma\tau(Z)=ZC_\sigma C_\tau\)，所以这是群同态。\(Z\) 的系数生成 \(L\)，故核平凡。若 \(Z'\) 是另一基本解矩阵，同一求导计算给出 \(Z'=ZB\)、\(B\in\operatorname{GL}_n(C)\)，于是新矩阵为 \(B^{-1}C_\sigma B\)。
\end{proof}

\begin{corollary}\label{b5:cor:E-group-closed}
设 \((K,\delta)\) 为特征零微分域，常数域 \(C=\ker\delta\) 代数闭，\(n\ge1\)，\(A\in M_n(K)\)，\(L/K\) 为系统 \(\delta(Y)=AY\) 的 PV 扩张，\(Z\in\operatorname{GL}_n(L)\) 为基本解矩阵。由 \(\sigma(Z)=ZC_\sigma\) 给出的微分 Galois 群矩阵像，是 \(\operatorname{GL}_n(C)\) 中的 Zariski 闭子群。
\end{corollary}
\begin{proof}
将 PV 环写成 \(R=K[X_{ij},(\det X)^{-1}]/I\)，其中 \(I\) 是极大真微分理想。取有限个理想生成元 \(f_1,\ldots,f_r\)。常数矩阵 \(c\) 给出自同构，恰当各个 \(f_j(Zc)=0\)：此条件让 \(X\mapsto Zc\) 通过商环；所得微分同态的像含 \(Z=(Zc)c^{-1}\)，所以满射，微分单性又保证单射。

对每个 \(j\)，把 \(f_j(Zc)\) 展开为 \(R\) 中有限多个 \(C\) 线性无关元素的组合，其系数是 \(c_{ij}\) 与 \((\det c)^{-1}\) 的多项式。它为零恰当这些有限个系数全部为零。这是一组 \(\operatorname{GL}_n(C)\) 上的正则方程，因而给出闭性；其含义正是保持解矩阵的全部代数关系。亦见\cite[Proposition 1.3.6, p. 12]{Singer2009DifferentialGalois}。
\end{proof}

下一节会直接检查具体关系及全部允许的变换。

\subsection{中间微分域与 Zariski 闭子群}
\label{b5:subsec:E02:03}

\begin{theorem}[微分 Galois 对应]\label{b5:thm:E-galois-correspondence}
设 \((K,\delta)\) 为特征零微分域，常数域 \(C=\ker\delta\) 代数闭，\(L/K\) 为 PV 扩张，\(G=\operatorname{Gal}_\delta(L/K)\) 为相应的 \(C\) 上线性代数群。则
\[
 H\longmapsto L^H,\qquad
 F\longmapsto\operatorname{Gal}_\delta(L/F)
\]
给出 \(G\) 的 Zariski 闭子群与中间微分域 \(K\subset F\subset L\) 之间逆序的双射。其中 \(L^H\) 是被 \(H\) 全部元素固定的子域。对应的 \(F/K\) 为 PV 扩张当且仅当 \(H\) 正规，此时 \(\operatorname{Gal}_\delta(F/K)\cong G/H\)。
\end{theorem}

\begin{proofsketch}
记 \(R\) 为 PV 环。在 \(R\otimes_KR\) 中，两基本解矩阵
\(Z_1=Z\otimes1\)、\(Z_2=1\otimes Z\) 的比较矩阵
\(U=Z_1^{-1}Z_2\) 满足 \(U'=0\)。证明\cref{b5:thm:E-pv-existence}时的常数变量引理，先在 \(L\otimes_KR\) 中把全部关系化为常数关系，再收缩回 \(R\otimes_KR\)，得到
\[
 R\otimes_KR\simeq R\otimes_C C[G],\qquad Z_2=Z_1U.
\]
这些常数关系的零点 \(c\) 给出 \(Z\mapsto Zc\) 的微分自同构，反向每个微分自同构也给出一个零点。特征零使这里的张量代数约简，故常数关系的理想是根理想，由 Nullstellensatz 它正好定义闭群 \(G\)，商环即 \(C[G]\)。这个同构表示 \(\operatorname{Spec}R\) 是一个 \(G\) 主齐性空间：选定一组解后，另一组解由唯一的群元素联系。

若 \(f\in L\) 被全部 \(G\) 元素固定，群常数点的 Zariski 稠密性使上述比较恒等式给出
\(f\otimes1=1\otimes f\) 在 \(L\otimes_KL\) 中成立。域扩张的忠实平坦下降于是推出 \(f\in K\)。该最后步骤也可取一组含一的 \(K\) 基直接比较张量系数。对任意中间微分域 \(F\)，\(L\) 仍由同一 \(Z\) 在 \(F\) 上生成，且常数不变；\cref{b5:thm:E-pv-existence}说明它是 \(F\) 上的 PV 域。重复这一步得到
\(L^{\operatorname{Gal}_\delta(L/F)}=F\)，而固定 \(F\) 的群在保持代数关系的方程中增加固定各个 \(f\in F\) 的方程，仍为 Zariski 闭群。

另一个方向使用闭子群的有理不变量分离陪集引理：线性代数群 \(G\) 的闭子群 \(H\) 有齐性空间 \(G/H\)，其有理函数能分离一般陪集，故固定全部右 \(H\) 不变有理函数的右平移恰来自 \(H\)。构造这个齐性空间时，取定义 \(H\) 的理想的一个有限维平移稳定生成空间，将其相应子空间或外幂直线嵌入 Grassmann 簇，\(H\) 正是该点的稳定子群。对上述主齐性空间作同样的下降，就得到
\(\operatorname{Gal}_\delta(L/L^H)=H\)。这里省略齐性空间的代数构造和有理函数的下降验证；它们是这一方向的主要技术步骤。

当 \(H\) 正规时，\(G/H\) 是线性代数群，主齐性空间的商可用它的一组忠实矩阵坐标描述。把这些坐标写成解矩阵的有限张量构造，导子给出一个 \(K\) 系数线性系统，其解坐标生成 \(L^H\)，所以该域是 PV 域。反之，若 \(F/K\) 是 PV 扩张，\(G\) 的每个元素都把它的基本解矩阵变成同一矩阵乘一个可逆常数矩阵，故保持 \(F\)。限制自同构因此得到 \(G\to\operatorname{Gal}_\delta(F/K)\)。\(F\) 的任意微分 \(K\) 自同构 \(\tau\) 都能延伸到 \(L\)：将 \(L/F\) 的系统沿 \(\tau\) 搬运，因其系数在 \(K\) 中而得到同一系统，再用 PV 唯一性取得延伸。因此限制满射，核为 \(H\)，故 \(H\) 正规且商群为所述 Galois 群。这部分省略商主齐性空间的坐标与线性系统的下降。详细的主齐性空间及对应证明见\cite[Section 1.4, Theorem 1.28, Proposition 1.34 and Corollary 1.40]{VanDerPutSinger2003GaloisTheory}，结果的概述见\cite[Theorem 1.3.9 and Proposition 1.3.20, pp. 13--17]{Singer2009DifferentialGalois}。
\end{proofsketch}

对应中的中间域须对导子稳定，子群须为 Zariski 闭子群，两项条件都不可删去。

\ProseParagraphBreak
例如在下一节的乘法群例子中，子群 \(\{2^m:m\in\mathbb Z\}\subset\mathbb C^\times\) 是真子群，却 Zariski 稠密：一元非零多项式不可能有无限多个根。若有理函数 \(f(t)\in K(t)\) 被这些伸缩全部固定，恒等式 \(f(ct)=f(t)\) 清去分母后是关于 \(c\) 的多项式等式，既然在无限多个 \(c\) 上成立，就在所有 \(c\ne0\) 上成立。故该真子群与整个乘法群有相同不动域，不能在删去闭性后仍要求双射。
